\section{Introduction}
\label{sec:intro}

\subsection{The question}

This paper studies a single fermion field with sixteen complex components on a spacetime of eight
dimensions. Two such spacetimes are of interest. One has four space-like and four time-like
dimensions, and is written $4+4$; the other has seven space-like dimensions and one time-like
dimension, and is written $7+1$. On both, the eight Dirac matrices can be chosen real
\cite{Nash1984,Nash1990,Nash2010}, and the field
supports a superfluid of fermion pairs of spin zero whose order parameter points along one of the
four coordinates that are not the coordinates of ordinary spacetime. \SIG{This paper treats the $4+4$
spacetime. Spacetimes with more than one time have been studied for their own sake
\cite{Bars2001}; what a second time does to the evolution of a field is taken up in
Sec.~\ref{sec:modes}.}{This paper treats the $7+1$ spacetime.}

The eight coordinates are written $x0,x1,\dots,x7$. Throughout, $x1,x2,x3$ are the coordinates of
ordinary space and $x4$ is the ordinary time. The coordinate $x0$ is an additional space-like
coordinate, called the hidden coordinate. The remaining three, $x5,x6,x7$, are
\SIG{three additional \emph{time-like} coordinates}{three additional \emph{space-like} coordinates}.

On a flat spacetime nothing explains why the fields of ordinary physics do not depend on $x5$, $x6$
and $x7$. \SIG{On the $4+4$ spacetime this is more than a gap in the explanation. A field that
depends on an additional time does not evolve unitarily in the ordinary time, and on flat spacetime
one has to \emph{impose} that the fields do not depend on the additional times.}{On a flat $7+1$
spacetime it has to be assumed.} The present paper examines a definite proposal for what happens to
these three coordinates, and it does so for a field with a mass and with a self-interaction.

\subsection{The hypothesis of deflation}

The hypothesis is that \emph{the coordinates $x5$, $x6$ and $x7$ deflate exponentially}: as the
ordinary time $x4$ advances, the proper extent of every interval along each of these coordinates
shrinks as an exponential of $x4$. Flat spacetime is replaced by a curved spacetime with a diagonal
metric, in which the three components of the metric along the deflating coordinates are
\begin{equation}
\begin{aligned}
\gm_{55}&=\SIG{-}{+}\cff^{2}\,e^{-2\,\gff\,x4},\\
\gm_{66}&=\SIG{-}{+}\css^{2}\,e^{-2\,\gss\,x4},\\
\gm_{77}&=\SIG{-}{+}\cvv^{2}\,e^{-2\,\gvv\,x4}.
\end{aligned}
\label{eq:hyp}
\end{equation}
Here $\gff$, $\gss$ and $\gvv$ are three positive constants, the \emph{rates} of deflation; $\cff$,
$\css$ and $\cvv$ are three positive constants that set the scale at $x4=0$ and may differ from one
another. The names of the rates are written upright and without a subscript, to distinguish them
from the components $\gm_{55}$, $\gm_{66}$, $\gm_{77}$ of the metric, which carry subscripts and the
superscript \SIG{$(4+4)$}{$(7+1)$}. The other five components of the metric keep their flat values,
$+1$ for $x0,\dots,x3$ and $-1$ for the time $x4$. The sign convention is that space-like components
of the metric are positive and time-like components negative: \SIG{the minus signs in
Eq.~(\ref{eq:hyp}) say that $x5,x6,x7$ are time-like.}{the plus signs in Eq.~(\ref{eq:hyp}) say that
$x5,x6,x7$ are space-like.}

\subsection{Mass term, self-interaction, and the two methods}

The Lagrangian of the fermion field has three parts (Sec.~\ref{sec:dirac}): the kinetic term; a
\emph{mass term}, bilinear in the field; and a \emph{self-interaction}, quartic in the field. The
self-interaction is the most general one without derivatives: nine invariants, one for each number
$r=0,\dots,8$ of Dirac matrices in the bilinears from which it is built, with nine couplings $g_r$.
The field equation is a nonlinear Dirac equation. The paper answers the following questions about
it. How many of the nine couplings are independent? What do the mass term and the self-interaction
do to matter made of these fermions, on flat spacetime and while three coordinates deflate? What do
they do to the pairing of the fermions, which in this model is mediated by a second field, the
Vorton-Scalar? And what does such matter do to the spacetime itself, if it is the source of gravity?

Two methods are used throughout, and each checks the other.
\begin{itemize}
\item The \emph{Kohn-Sham method}, known from density-functional theory: the interacting fermions
are replaced by independent fermions that move in fields which are determined self-consistently by
their own densities. For uniform matter the fields are a shift of the mass and a shift of the
energy.
\item The \emph{effective potential for composite operators} of Cornwall, Jackiw and Tomboulis
(CJT): a function of the full propagator of the fermions whose stationary points are the possible
states of matter and whose value there is minus the pressure.
\end{itemize}
The two methods are equivalent at the level of approximation used here, and the paper verifies this
numerically (Sec.~\ref{sec:cjtcontact}).

The paper also examines three hypotheses about the superfluid.
\begin{itemize}
\item[H1.] The fermion field forms a superfluid state of spin 0 in which the fermions are paired by
the exchange of a new scalar boson, the \emph{Vorton-Scalar}, through a neutral current of the kind
known from the weak interaction.
\item[H2.] This superfluid state is the Higgs field.
\item[H3.] This superfluid state explains the observations of dark matter, or the Hubble tension, or
both.
\end{itemize}
The purpose is to carry out every calculation that these hypotheses call for under the hypothesis of
deflation, with the mass term and the self-interaction, and to report what follows for H1, H2 and
H3, favorable or not.

\subsection{Summary of results}

\emph{Geometry} (Sec.~\ref{sec:geometry}). The deflating spacetime is homogeneous and has 14
symmetries, or 15 or 20 if two or three of the rates are equal, instead of the 36 of flat spacetime.
It is not a solution of Einstein's equations without matter. With three equal rates it \emph{is} an
exact solution of the equations of Einstein-Gauss-Bonnet gravity without matter, for one definite
relation between the rate and the two constants of that theory.

\emph{The Dirac field} (Secs.~\ref{sec:dirac}, \ref{sec:modes}). A field that does not depend on
$x5,x6,x7$ obeys, after multiplication by the square root of the volume factor of the deflating
coordinates, exactly the flat equation with the mass unchanged and with every coupling of the
self-interaction growing as the exponential of the time. A field that does depend on the deflating
coordinates is affected directly.
%%<1644>
The linear equation is solved exactly. Of its sixteen solutions eight grow as the exponential of an
exponential of the time and are not admissible. The other eight are standing waves, which are
extinguished as the exponential of an exponential. \emph{The deflation removes the dependence of the
field on the additional times}, and with it the sector in which the $4+4$ theory is not unitary.
%%</1644>
%%<1671>
The linear equation is solved exactly. \emph{The deflation creates particles}: a mode that starts
empty ends with the occupation of a Fermi-Dirac distribution at a temperature equal to the rate
divided by $2\pi$, the Gibbons-Hawking temperature of the de Sitter space that the deflating
coordinates and the time form when the rates are equal.
%%</1671>

\emph{The self-interaction} (Sec.~\ref{sec:selfint}). For anticommuting fields the nine invariants
obey five linear relations: only four couplings are independent, a scalar, a vector, an axial-vector
and a pseudoscalar one. Uniform matter feels two combinations of them, pairs of particles at rest
feel a third. \SIG{On the $4+4$ spacetime the self-interaction has two forms in the sector that can
be quantized, which agree for particles in the same half of the field and differ in sign for
particles in different halves. The covariant form is the definition of the model in this paper;
the other is recorded beside it, and the two are discussed in Sec.~\ref{sec:formsdisc}.}{}

\emph{Matter} (Secs.~\ref{sec:ks}, \ref{sec:cjtcontact}). The Kohn-Sham method and the CJT effective
potential give the same states. \SIG{}{For matter that has momentum along all seven space coordinates
the deflation raises the density and the anisotropy of the pressure, and matter that is self-bound
at one moment is so only for about one e-folding of the deflating volume. }In the sector of fields
that do not depend on the deflating coordinates a scalar attraction grows with the deflation and
turns a gas into self-bound matter at a definite moment. With the filled states of negative energy
taken into account, the same growth gives a massless field a mass at a critical coupling.\SIG{ In
the covariant form matter polarizes: the self-bound state fills one half of the field and forms at
two to three times the strength that the alternative form needs, and the mass that a massless field
acquires has opposite signs in the two halves.}{} A bare
mass together with the Vorton-Scalar gives the phase mode of that field a mass.

\emph{The superfluid} (Secs.~\ref{sec:clock} to \ref{sec:order}). For fields that do not depend on
the extra coordinates the pairing calculation of flat spacetime holds at each moment with the
couplings of that moment. The masses do not change. Every threshold is passed at a definite time. The
gap of the paired state grows much faster than the coupling; \SIG{in the pairs along the additional
times a scalar contact attraction raises it and a vector contact repulsion lowers it, and in the
pair along the hidden axis the covariant form reverses both}{a scalar contact attraction raises it,
a vector contact repulsion lowers it}. The gap equation has a solution that does not depend on the
upper limit of its momentum integral only below an ultraviolet criterion. Two fermions become bound at a scattering resonance. The
paired state can follow the deflation only once its gap exceeds a few times the rate divided by the
pairing strength; before that it lags behind and breaks pairs. Matter of a given density loses its
fermion mass at a definite time. A vortex line carries four fermion zero modes.

\emph{Gravity} (Sec.~\ref{sec:gravity}). With Kohn-Sham matter as the source of the Einstein-Lovelock
equations the deflating spacetime is not stationary: the three coordinates deflate faster than
exponentially and the four flat coordinates expand. Even without matter the exactly exponential
deflation is an unstable solution of the gravity in which it is a solution.

\emph{Cosmology} (Sec.~\ref{sec:cosmo}). If gravity is the Einstein theory in eight dimensions,
Newton's constant grows at the total rate of deflation; lunar laser ranging then bounds that rate to
$3\times10^{-3}$ of the Hubble rate, and the Hubble constant inferred from the microwave background
falls. The deflation does not relieve the Hubble tension. Whatever keeps the rates constant must
carry a tension and must take part in the perturbations of ordinary space; the spectra of the
microwave background at all but the largest angular scales do not depend on how it does so, and
the growth of structure does.

\emph{The Higgs hypothesis} (Sec.~\ref{sec:higgs}). \SIG{The symmetry under which the pair condensate
has the pattern of the Higgs vacuum is a hidden symmetry on the $4+4$ spacetime, flat or deflating.}{The
rotations of the four extra space-like coordinates, under which the pair condensate has the pattern
of the Higgs vacuum, are no longer symmetries of the spacetime; the pattern survives only as a hidden
symmetry.} \SIG{A contact self-interaction of the covariant form breaks it explicitly, down to the
rotations of the additional times.}{A vector or an axial-vector contact coupling breaks one factor of
that symmetry.} If the
condensate were the Higgs field, the electroweak scale would follow the growing coupling, and atomic
clocks would limit the rate of deflation to about $10^{-6}$ of the Hubble rate. The hypothesis H2 and
a deflation of any consequence today exclude each other.

\subsection{Method, and the limits of what is claimed}

Every statement of this paper about geometry, matrices and operators is an identity that a program
has verified in exact rational arithmetic, and that an independent program has verified again by
computer algebra or with Grassmann-valued fields. Every numerical statement is the output of a
program, with a stated accuracy, and the central ones were computed twice by independent programs.
The programs and their output are public \cite{DeflationRepo}; Appendix~\ref{app:checks} says how each
number of this paper is reproduced. The program performs {{CHECKS}} checks and the independent
verification {{GEN:n_indep_checks}}, all of which pass. The metric of Eq.~(\ref{eq:hyp}) is a
hypothesis. The paper derives its consequences; it does not claim that the hypothesis is true.
Section~\ref{sec:status} lists what has been established and what has not.

Units are natural, $\hbar=c=k_B=1$. Symbols that carry the tag $\tg$ refer to the deflating
\SIG{$4+4$}{$7+1$} spacetime. A general coordinate among the three deflating ones is written $xt$ with
$t\in\{5,6,7\}$ and its rate $\rt$; a general coordinate among $x0,x1,x2,x3$ is written $xj$ with
$j\in\{0,1,2,3\}$. The letter $a$ is kept for the four extra axes, $a\in\{0,5,6,7\}$
(Sec.~\ref{sec:gap}). Every symbol is defined where it first appears, and Appendix~\ref{app:symbols}
lists all of them.
